\chapter{Cached Context}
\label{cha:ch3}

\myx{Please ping me when you are done with this chapter.}

\section{Motivation}
\label{sec:ch3_intro}

As we have identified the existence of \textbf{Challenge 1} (poor patch quality due to incomplete context) and \textbf{Challenge 2} (slow generation due to long context) in modern agentic systems, it is therefore natural to consider designing an end-to-end system that aims at tackling both challenges, given the combined performance, latency, and cost benefits of a joint optimization. In the previous chapters, we evaluated the performance of modern solutions for each respective challenge, which are RepoGraph \cite{ouyang2025repographenhancingaisoftware} and CacheBlend \cite{yao2025cacheblendfastlargelanguage}. RepoGraph extracts repository-level context via pre-generating a structure tree containing symbol locations and a code graph containing call relations. CacheBlend enables position-independent KV cache full reuse by selectively recomputes a small subset of high-KV-deviation tokens to restore cross-attention. {\colorGreen In our integrated system, the aim is to combine the performance and latency advantages of each while mitigating the degradation of accuracy due to partially recomputed KV cache. The core concept that allows this is \textit{Cached Context}, which is a simple yet effective idea of chunking context blocks in prompts generated in the agentic workflow, such that the back-end KV cache reuse scheme (CacheBlend in our case) is able to exploit reuse opportunities more frequently. The reason why this is particularly effective in agentic settings is primarily because of the modular and repetitive nature of code agent context, which mostly consists of repository structure related info and function or class definition code blocks. Prompt templates, often containing format requirements and task-specific instructions, also will have tokens that are repeated upon every new request sent to the LLM. \textit{Cached Context} allows capturing these additional reuse opportunities to further speed up inference without loss of accuracy.}